\documentclass[a4paper,10pt]{article}
\usepackage[utf8]{inputenc}
\usepackage[T1]{fontenc}
\usepackage{lmodern}
\usepackage[margin=1.4cm, top=1.2cm, bottom=1.2cm]{geometry}
\usepackage{hyperref}
\usepackage{xcolor}
\usepackage{enumitem}
\usepackage{graphicx}
\usepackage{titlesec}

% Minimal clean styling
\hypersetup{colorlinks=true, urlcolor=blue!70!black, linkcolor=blue!70!black}
\pagestyle{empty}
\setlength{\parindent}{0pt}

% Section formatting
\titleformat{\section}{\large\bfseries\color{black}}{}{0em}{}[\titlerule]
\titlespacing*{\section}{0pt}{6pt}{4pt}

% Colors
\definecolor{accent}{RGB}{20, 50, 120}

% Role command
\newcommand{\role}[4]{%
  \vspace{2pt}\noindent
  \begin{minipage}[t]{0.70\textwidth}\raggedright\textbf{#1} \textbullet\ \textit{#3}\end{minipage}%
  \hfill
  \begin{minipage}[t]{0.28\textwidth}\raggedleft\footnotesize\color{gray}#2 \textbullet\ #4\end{minipage}%
  \vspace{1pt}%
}
\setlist[itemize]{leftmargin=1.2em, itemsep=1.5pt, topsep=2pt, parsep=0pt}

\begin{document}

% ================= HEADER =================
\begin{minipage}[c]{0.80\textwidth}
  {\Huge\bfseries Prabhavit Gupta}\\[3pt]
  {\color{accent}\large ML $\&$ AI Engineer $\vert$ MLOps, Generative AI, AI Agents, Databricks $\&$ Cloud AI}\\[6pt]
  {\footnotesize +39 351 610 7807 ~$\vert$~
  \href{mailto:prabhavitg@gmail.com}{prabhavitg@gmail.com} ~$\vert$~
  \href{https://www.linkedin.com/in/prabhavit/}{linkedin.com/in/prabhavit} ~$\vert$~
  \href{https://github.com/prab-gupta}{github.com/prab-gupta} ~$\vert$~
  Venezia / Milano, Italia (Residente Permanente UE $\vert$ Full Remote $\vert$ Disponibilità Immediata)}
\end{minipage}\hfill

% ================= SUMMARY =================
\section{Summary}
ML e AI Engineer con 3 anni di esperienza commerciale nella progettazione, sviluppo e industrializzazione end-to-end di soluzioni di Machine Learning, Generative AI e sistemi agentici. Esperto nella realizzazione di pipeline MLOps per la gestione completa del ciclo di vita dei modelli (MLflow, CI/CD, Docker, Kubernetes) e nell'integrazione di servizi LLM aziendali (Azure OpenAI, Gemini, Anthropic) con architetture RAG scalabili su Databricks e PySpark. Comprovata competenza nella definizione di architetture basate su AI Agents, governance, sicurezza e explainability per contesti enterprise ad alto impatto. Residente permanente a Venezia, immediatamente disponibile per ingaggi full remote / ibridi su Milano.

% ================= EXPERIENCE =================
\section{Experience}

\role{\href{https://flowspace.co}{FlowSpace.co}}{Mar 2024 - Aug 2026}{Lead AI Engineer $\&$ MLOps Architect (Freelance)}{Londra, UK (remoto)}
\begin{itemize}
  \item Guidato lo sviluppo e l'industrializzazione di modelli ML e soluzioni Generative AI in produzione per oltre 4.500 utenti enterprise attivi.
  \item Progettato e implementato pipeline MLOps complete (tracking esperimenti, model registry con MLflow, containerizzazione Docker e CI/CD) per la gestione affidabile del ciclo di vita dei modelli.
  \item Architettato sistemi multi-agente autonomi (AI Agents con LangGraph) e flussi RAG avanzati integrati con vector store (pgvector, Chroma) e servizi cloud (Azure, GCP, AWS Bedrock).
  \item Integrato workflow analitici distribuiti su Databricks per l'elaborazione su larga scala di dataset testuali e la generazione di metriche di validazione automatizzate.
  \item Implementato best practice per la governance, mitigazione di allucinazioni, sicurezza e monitoraggio continuo di latenza (<1.8s) e drift prestazionale in esercizio.
  \item \textit{Tecnologie: Python, PyTorch, Generative AI, MLOps, MLflow, Databricks, Azure OpenAI, Docker, Kubernetes, FastAPI, Git}
\end{itemize}

\role{Urlink}{Apr 2025 - Jul 2025}{Machine Learning $\&$ Data Engineer (Freelance)}{Barcellona, Spagna (remoto)}
\begin{itemize}
  \item Sviluppato pipeline per l'ingestion, normalizzazione semantica e classificazione automatica di flussi informativi non strutturati.
  \item Realizzato microservizi in Python e FastAPI con monitoraggio automatico della qualità dei dati e validazione schemi Pydantic.
  \item Collaborato con team di data platform per ottimizzare query SQL e performance di inferenza in staging e produzione.
  \item \textit{Tecnologie: Python, Scikit-learn, PyTorch, Pandas, PostgreSQL, Docker, REST APIs, Git}
\end{itemize}

\role{Consulente Tecnico Indipendente}{Sep 2023 - Mar 2024}{AI Solutions Engineer}{Venezia, Italia}
\begin{itemize}
  \item Progettato e realizzato soluzioni su misura di Generative AI, assistenti conversazionali e pipeline RAG per aziende italiane ed europee.
  \item Guidato i clienti nell'adozione di architetture agentiche per l'automazione di workflow documentali e gestionali complessi.
  \item Redatto specifiche architetturali, documentazione tecnica di governance e runbook operativi per il passaggio in produzione.
  \item \textit{Tecnologie: Python, Azure AI, LangChain, Vector DB, FastAPI, Linux, Git}
\end{itemize}

\role{\href{https://tractionboard.co}{Tractionboard.co}}{Aug 2023 - Oct 2023}{Full-Stack Automation Developer (Stage)}{Utrecht, Paesi Bassi (remoto)}
\begin{itemize}
  \item Sviluppato script di data ingestion e connettori webhook per piattaforme SaaS di monitoraggio analytics.
  \item \textit{Tecnologie: Python, JavaScript, REST APIs, Git}
\end{itemize}

% ================= EDUCATION =================
\section{Education}
\role{Università Ca' Foscari di Venezia}{Ott 2020 - Giu 2024}{Laurea Triennale in Filosofia, Studi Internazionali ed Economia, Minor in Imprenditorialità}{Venezia, Italia}

% ================= CERTIFICATIONS =================
\section{Certifications}
Microsoft: Introduction to building AI agents \textbullet\ Microsoft Copilot Studio Agent Academy \textbullet\ Microsoft Power Automate \textbullet\ Tray.AI Build Practitioner (Sep 2024) \textbullet\ Tray.AI Build Foundations

% ================= ADDITIONAL INFORMATION =================
\section{Additional Information}
\begin{itemize}
  \item \textbf{Inquadramento:} Residente permanente UE a Venezia, Italia. Disponibile per contratti CCNL o collaborazione con Partita IVA attiva.
  \item \textbf{Lingue:} Italiano (lavorativo / professionale), Inglese (fluente C1 professionale), Hindi (madrelingua).
  \item \textbf{Disponibilità:} Immediata (0 giorni di preavviso) per posizioni full remote / ibride su Milano.
\end{itemize}

% ================= CORE SKILLS =================
\section{Core Skills}
\begin{itemize}
  \item \textbf{Machine Learning $\&$ MLOps:} Model Training $\&$ Deployment, MLflow, Model Registry, Monitoring in Esercizio, PyTorch, Scikit-learn, PySpark.
  \item \textbf{Generative AI $\&$ AI Agents:} LLM (Azure OpenAI, Gemini, Claude, Llama), Architetture Multi-Agent (LangGraph), Prompt Engineering, RAG Avanzato.
  \item \textbf{Piattaforme Enterprise $\&$ Cloud:} Databricks, Apache Spark, Azure, Google Cloud, AWS, Kubernetes, Docker, Microservizi FastAPI.
  \item \textbf{Governance $\&$ Software Engineering:} Python 3.12, Clean Code, CI/CD, Sicurezza, Explainability, Git, Linux.
\end{itemize}

\end{document}
